\documentclass[11pt,a4paper]{article}
\usepackage[margin=2.2cm]{geometry}
\usepackage[T1]{fontenc}
\usepackage{lmodern}
\usepackage{microtype}
\usepackage{amsmath}
\usepackage{booktabs}
\usepackage{xcolor}
\usepackage{listings}
\usepackage[most]{tcolorbox}
\usepackage[colorlinks=true,linkcolor=teal!70!black,urlcolor=teal!70!black]{hyperref}
\usepackage{enumitem}
\setlist{itemsep=2pt,topsep=4pt}

\definecolor{codebg}{RGB}{246,247,249}
\definecolor{kw}{RGB}{29,106,134}
\definecolor{cm}{RGB}{110,120,130}
\definecolor{st}{RGB}{176,48,40}
\lstdefinestyle{m}{language=Matlab,basicstyle=\ttfamily\small,keywordstyle=\color{kw}\bfseries,
  commentstyle=\color{cm}\itshape,stringstyle=\color{st},backgroundcolor=\color{codebg},
  frame=single,rulecolor=\color{black!12},numbers=left,numberstyle=\tiny\color{cm},
  xleftmargin=1.6em,framexleftmargin=1.4em,breaklines=true,columns=fullflexible,keepspaces=true,
  morekeywords={daq,addinput,read,seconds,save,saveas,struct,disp,fprintf,num2str,figure,plot,xlabel,ylabel,title,grid,isfile,error,datetime,input}}
\lstdefinestyle{py}{language=Python,basicstyle=\ttfamily\small,keywordstyle=\color{kw}\bfseries,
  commentstyle=\color{cm}\itshape,stringstyle=\color{st},backgroundcolor=\color{codebg},
  frame=single,rulecolor=\color{black!12},breaklines=true,columns=fullflexible,keepspaces=true}
\newtcolorbox{keypoint}[1][]{colback=teal!5,colframe=teal!50!black,boxrule=0.5pt,arc=2pt,left=6pt,right=6pt,top=4pt,bottom=4pt,fonttitle=\bfseries\small,title={#1}}
\newtcolorbox{warn}[1][]{colback=orange!6,colframe=orange!60!black,boxrule=0.5pt,arc=2pt,left=6pt,right=6pt,top=4pt,bottom=4pt,fonttitle=\bfseries\small,title={#1}}
\newcommand{\code}[1]{\texttt{#1}}

\title{\vspace{-1.2cm}\textbf{The FLDI acquisition script, explained}\\[2pt]
\large \code{DAQ\_FLDI\_Setup.m}, summer 2025 hot-air-gun campaign}
\author{}
\date{}

\begin{document}
\maketitle
\vspace{-1.6cm}

\section{What the script does, in one paragraph}
Each time you pressed \emph{Run}, the script recorded the photodiode voltage for one test point. It
connects to a National Instruments (NI) data-acquisition card, records 5\,s of voltage from one
input channel at 10\,000 samples per second, computes the size of the fluctuations (their RMS),
saves everything to a file named \code{t<run number>.mat}, and saves a picture of the voltage
trace as \code{t<run number>.jpg}. Before each run you edited two lines by hand: the run number
and the position label. Nothing else changed from run to run.

\section{The full script}
\begin{lstlisting}[style=m]
close all
clear
clc
format longEng

%%
RNo = 133;
loc = '30 mm (diode)';
deviceID = 'Dev2';
channelID = 'ai0';
fs = 10000;           % Sampling frequency in Hz
numSamples = 50000;   % Number of samples to acquire
duration = numSamples / fs;

% Create DAQ session
dq = daq("ni");
ch = addinput(dq, deviceID, channelID, "Voltage");
ch.Range = [-1 1];

% Configure session
dq.Rate = fs;
% dq.DurationInSeconds = duration;

% Display acquisition start
disp('Starting acquisition...');

% Acquire data
[data, timestamps] = read(dq, seconds(duration), "OutputFormat", "Matrix");

% Display acquisition complete
disp('Acquisition complete.');

% Save data with metadata
outputFile = 'daq_data.mat';
channelInfo = struct( ...
    'Device', deviceID, ...
    'Channel', channelID, ...
    'Rate', fs, ...
    'NumSamples', numSamples ...
);

% Compute RMS after subtracting mean
meanSubtracted = data - mean(data);
rmsValue = sqrt(mean(meanSubtracted.^2));

disp('Writing start.');
filenameD = ['t' num2str(RNo) '.mat'];
save(filenameD, 'data', 'timestamps', 'channelInfo', 'loc', 'fs', 'duration', 'numSamples', 'RNo', 'rmsValue');
disp('Writing end.');

% Display RMS
fprintf('Mean-subtracted RMS voltage: %.4f V\n', rmsValue);

% Plot the acquired data
timeVector = (0:numSamples-1) / fs;  % Time in seconds
figure;
plot(timeVector, data);
xlabel('Time (s)');
ylabel('Voltage (V)');
title('Acquired Voltage Signal');
grid on;
filename = ['t' num2str(RNo) '.jpg'];
saveas(gcf, filename);
\end{lstlisting}

\section{Block by block}

\subsection*{Lines 1--4: start from a clean slate}
\code{close all} closes old figure windows, \code{clear} deletes all variables left over from the
previous run, and \code{clc} clears the command window. \code{format longEng} only changes how
numbers are printed on screen (engineering notation, e.g.\ \code{12.3456e-003}); it has no effect
on the data. Clearing variables matters: without it, a value from the previous run could silently
be reused.

\subsection*{Lines 7--13: the settings for this run}
\begin{itemize}
\item \code{RNo} is the run number. It becomes the file name, so \textbf{two runs with the same
  number overwrite each other}.
\item \code{loc} is a free-text label for where the hot-air gun was, e.g.\ \code{'30 mm (diode)'}
  means 30\,mm from the focus towards the photodiode; \code{(laser)} means towards the laser.
  Other labels used: \code{'heat gun off'}, \code{'wollaston test'}, \code{'wollaston test V0'}.
\item \code{deviceID = 'Dev2'} and \code{channelID = 'ai0'} pick the NI card and its analogue
  input 0, where the photodiode cable was plugged in. ``Dev2'' is simply the name NI's software
  gave the card on that computer.
\item \code{fs = 10000} is the sampling rate (samples per second); \code{numSamples = 50000} is
  how many samples to take. So \code{duration} $= 50000/10000 = 5$\,s.
\end{itemize}

\subsection*{Lines 16--22: configure the DAQ card}
\code{daq("ni")} opens a connection to NI hardware through MATLAB's Data Acquisition Toolbox.
\code{addinput} adds one voltage input channel. \code{ch.Range = [-1 1]} tells the card to expect
voltages between $-1$ and $+1$\,V. A smaller range gives finer voltage steps; anything outside the
range is clipped. The photodiode signal sat between about 0.01 and 0.39\,V, so this range fitted
well. \code{dq.Rate = fs} sets 10\,kHz sampling. Line 22 is commented out and does nothing; the
duration is given directly to \code{read} on line 28 instead.

\subsection*{Line 28: record}
\code{read(dq, seconds(5), ...)} blocks for 5\,s while the card records, then returns two column
vectors with 50\,000 entries each: \code{data} (volts) and \code{timestamps} (seconds from the
start, 0, 0.0001, 0.0002, \dots).

\subsection*{Lines 34--49: save}
\code{channelInfo} bundles the hardware settings into one structure so they travel with the data.
\code{outputFile = 'daq\_data.mat'} is never used (leftover code). The file name is built as
\code{'t'} + run number + \code{'.mat'}, and \code{save} writes nine variables into it (listed in
Section~5).

\subsection*{Lines 43--44 and 52: the RMS value}
The script subtracts the average voltage and takes the root-mean-square of what remains:
\[
V_{\mathrm{rms}} = \sqrt{\frac{1}{N}\sum_{i=1}^{N}\big(V_i-\bar V\big)^2}.
\]
This is the standard deviation of the signal, i.e.\ the typical size of the fluctuations. The mean
voltage $\bar V$ is the steady light level, set by where the interferometer is biased on its
fringe; the fluctuations are what the flow causes. This printed number is the ``V'' column you
copied into the Excel log after each run.

\subsection*{Lines 55--63: quick-look plot}
Plots voltage against time and saves it as a JPEG, so you could check by eye that the run looked
sensible (no clipping, no dropout). The images ended up in a folder called \code{PNGs}, but they
are JPEG files.

\section{What the settings imply}
\begin{center}
\begin{tabular}{@{}lll@{}}
\toprule
Quantity & Value & Why it matters \\
\midrule
Sampling rate $f_s$ & 10\,kHz & one sample every 0.1\,ms \\
Record length & 5\,s, 50\,000 samples & longer gives more averaging \\
Highest resolvable frequency & $f_s/2 = 5$\,kHz & the Nyquist limit (see Section~6) \\
Finest spectral resolution & $1/5\,\mathrm{s} = 0.2$\,Hz & before averaging blocks \\
Input range & $\pm1$\,V & signal must stay inside \\
Voltage step & $\approx 31\,\mu$V & 16-bit card (measured from the data) \\
Gun-off noise level & 0.3--0.8\,mV RMS & the floor of every measurement \\
\bottomrule
\end{tabular}
\end{center}

\section{How the script was used in the experiment}

\paragraph{Run-number convention.} The Excel sheets show the same pattern every day:
\begin{itemize}
\item \textbf{Runs 1000--1002, ``Wollaston test''}: you slowly translated the Wollaston prism
  during the 5\,s record so the light swept through bright and dark fringes. The minimum and
  maximum voltages give $V_{\min}$ and $V_{\max}$.
\item \textbf{Run 1003 (or 1004), ``V0''}: the prism set back to the middle of the fringe and the
  steady voltage $V_0$ recorded. The middle, $V_0 \approx (V_{\min}+V_{\max})/2$, is the
  \emph{quadrature} point where the interferometer is most sensitive and most linear.
\item \textbf{Runs 1--6, ``heat gun off''}: the background noise of the instrument.
\item \textbf{Runs 7 onwards}: the gun at each traverse position, usually 2--3 repeats per
  position.
\end{itemize}

\paragraph{From the logged RMS to a physical quantity.} Near quadrature, the voltage changes by
$(V_{\max}-V_{\min})/2$ per radian of optical phase difference between the two FLDI beams. So the
RMS phase fluctuation caused by the hot air is
\[
\phi_{\mathrm{rms}} \approx \frac{2\,V_{\mathrm{rms}}}{V_{\max}-V_{\min}}.
\]
\begin{keypoint}[{Example: 12 June, gun at the focus (run 31)}]
$V_{\min}=0.0124$\,V, $V_{\max}=0.3759$\,V, $V_{\mathrm{rms}}=0.0338$\,V, giving
$\phi_{\mathrm{rms}} \approx 2(0.0338)/0.3635 = 0.19$\,rad. Converting to phase is what makes
different days comparable: on 3--4 June the fringe was about four times smaller
($V_{\max}-V_{\min}\approx0.095$\,V), so the same voltage meant four times more phase.
\end{keypoint}

\section{What is inside a saved file, and how to open it}
Each \code{t<RNo>.mat} contains: \code{data} (50\,000 voltages), \code{timestamps}, \code{fs},
\code{numSamples}, \code{duration}, \code{RNo}, \code{loc}, \code{rmsValue}, and
\code{channelInfo} (device, channel, rate, number of samples).

\noindent MATLAB, with a power spectrum:
\begin{lstlisting}[style=m,numbers=none]
S = load('t31.mat');
fprintf('%s, RMS = %.4f V\n', S.loc, S.rmsValue);
[P, f] = pwelch(S.data - mean(S.data), hann(4096), 2048, 4096, S.fs);
semilogy(f, P); xlabel('Frequency (Hz)'); ylabel('PSD (V^2/Hz)');
\end{lstlisting}
Python:
\begin{lstlisting}[style=py]
from scipy.io import loadmat
from scipy.signal import welch
S = loadmat('t31.mat')
v = S['data'].ravel(); fs = float(S['fs'])
print(S['loc'][0], v.std())          # same number as rmsValue
f, P = welch(v - v.mean(), fs=fs, nperseg=4096)
\end{lstlisting}

\section{Limitations to know about}
\begin{warn}[1. Nothing above 5\,kHz is measured correctly]
At 10\,kHz sampling, any signal faster than 5\,kHz is not just lost: it \emph{aliases}, i.e.\ it
appears at a false lower frequency. Most NI multifunction cards have no built-in filter to remove
it. In the 12 June data at the focus the spectrum is still high at 5\,kHz, so faster turbulence is
very likely folded into the 0--5\,kHz range and the RMS. The 30 June runs used 250\,kHz,
which removes this problem for flows like the hot-air gun.
\end{warn}
\begin{warn}[2. Run number and label are typed by hand]
A repeated \code{RNo} silently overwrites a file, and a forgotten \code{loc} edit mislabels it.
This happened at least once: on 4 June, \code{t1004.mat} carries the label \code{'0 mm (laser)'}
although it is a gun-off V0 check, and \code{t1003.mat} contains \code{RNo = 1004} inside.
Always trust the Excel log together with the file, and check any odd point.
\end{warn}
\begin{warn}[3. The RMS includes everything]
It contains instrument noise (0.3--0.8\,mV), slow drifts and all frequencies up to 5\,kHz.
For sensitivity studies it is better to compute the variance in a chosen frequency band from
\code{data}, rather than use \code{rmsValue}.
\end{warn}
\noindent Smaller points: the conditions of each run (gun temperature, date, time, distance from
nozzle to beam) are not stored in the file, only in the Excel sheet; $V_{\min}$, $V_{\max}$ are
not stored either; and \code{outputFile} and line~22 are unused.

\section{A safer version for future campaigns}
The changes below keep the same workflow but remove the manual-error risks and store what the
analysis needs. Set \code{fs} high enough for the flow (at least twice the highest frequency of
interest).
\begin{lstlisting}[style=m,numbers=none]
RNo   = 133;               % still set by hand
z_mm  = 30;                % position: + diode side, - laser side
note  = 'heat gun 200 C';  % free text
fs    = 250e3;  T = 5;     % rate and duration
Vmin  = 0.0124; Vmax = 0.3759;   % from today's Wollaston test

fname = sprintf('t%d.mat', RNo);
if isfile(fname), error('%s exists: change RNo', fname); end

dq = daq("ni"); ch = addinput(dq, "Dev2", "ai0", "Voltage");
ch.Range = [-1 1]; dq.Rate = fs;
[data, timestamps] = read(dq, seconds(T), "OutputFormat", "Matrix");

if max(abs(data)) > 0.98, warning('Signal close to range limit'); end
rmsValue = std(data, 1);
phiRms   = 2*rmsValue/(Vmax - Vmin);          % radians
when     = datetime('now');
save(fname, 'data','fs','T','RNo','z_mm','note','Vmin','Vmax', ...
     'rmsValue','phiRms','when');
fprintf('Run %d, z = %g mm: %.4f V rms, %.3f rad rms\n', RNo, z_mm, rmsValue, phiRms);
\end{lstlisting}

\end{document}
